\chapter{从变分原理到 Hartree--Fock}
\sourcepages{手写笔记第 18--30 页}
\section{求基态为什么可以变成最小化}
从本章起，电子公式默认使用 Hartree 原子单位：$\hbar=m_e=e=4\pi\epsilon_0=1$。长度单位为 $a_0=4\pi\epsilon_0\hbar^2/(m_e e^2)$，能量单位为 $E_h=\hbar^2/(m_ea_0^2)$。固定核时，先求不含核间常数的电子能量，最后另加 $E_{NN}$。

量子态可看成复向量，内积 $\langle\phi|\psi\rangle=\int\phi^*\psi$；位置连续时矩阵变成作用于函数的算符。归一化代表总概率为一。若 $\hat H|n\rangle=E_n|n\rangle$ 且 $E_0$ 最低，将试探态写为 $|\Psi\rangle=\sum_nc_n|n\rangle$，则
\begin{equation}
\langle\Psi|\hat H|\Psi\rangle-E_0=\sum_n|c_n|^2(E_n-E_0)\ge0.
\end{equation}
这就是变分原理。限制试探态的自由度，只能使所得最低能量高于或等于精确基态。等号要求试探态完全处于基态子空间。

泛函是把函数映为数的规则。若 $F[\phi+\delta\phi]-F[\phi]=\int (\delta F/\delta\phi)\delta\phi+O(\delta\phi^2)$，线性项系数就是泛函导数。复函数变分可将 $\phi,\phi^*$ 暂作独立变量。对 $\langle\phi|\hat H|\phi\rangle-\varepsilon(\langle\phi|\phi\rangle-1)$ 关于 $\phi^*$ 变分，便得 $\hat H\phi=\varepsilon\phi$。乘子确保最小化不能通过把波函数缩成零来作弊。

数值上可沿残差 $(\hat H-\varepsilon)\phi$ 的负方向移动，再归一化。归一化必须除以 $\sqrt{\langle\phi|\phi\rangle}$。求多个态时还要保持正交，否则所有向量都可能落向同一个最低态。非线性多轨道能量还可能存在多个驻点，自洽不等于已找到全局最低点。

\section{Hartree 乘积如何产生平均场}
多电子 Hamilton 算符为 $\hat H=\sum_i\hat h(i)+\sum_{i<j}r_{ij}^{-1}$，其中 $\hat h=-\nabla^2/2+v_{\rm ext}$。最简单的试探态是归一化单粒子函数的乘积 $\Psi_H=\prod_i\phi_i(x_i)$。$x=(\bm r,s)$ 包括位置与自旋，$\int\dd x$ 代表空间积分和自旋求和。

其他坐标的归一化积分等于一，因此一体项化为 $h_{ii}=\langle\phi_i|\hat h|\phi_i\rangle$；两体项为
\begin{equation}
J_{ij}=\iint\frac{|\phi_i(x)|^2|\phi_j(x')|^2}{|\bm r-\bm r'|}\dd x\dd x',\qquad
E_H=\sum_i h_{ii}+\frac12\sum_{i\ne j}J_{ij}.
\end{equation}
因每对粒子只相互作用一次，双重求和前有 $1/2$。对 $\phi_i^*$ 变分时，$J_{ij}$ 与 $J_{ji}$ 给出相同贡献，消去这个 $1/2$，得到
\begin{equation}
\left[\hat h+\sum_{j\ne i}\int\frac{|\phi_j(x')|^2}{|\bm r-\bm r'|}\dd x'\right]\phi_i=\varepsilon_i\phi_i.
\end{equation}
每个粒子看到其他粒子的平均密度。由于势依赖未知轨道，需要猜轨道、建势、解方程、更新轨道，直到输入与输出一致。这是自洽场方法。严格的 Hartree 乘积没有反对称性，不能作为电子的最终描述。

\section{反对称性怎样产生交换}
相同电子互换全部坐标时，波函数必须变号。用正交归一自旋轨道构造 Slater 行列式 $\Phi=(N!)^{-1/2}\det[\phi_i(x_j)]$，即可自动满足这一要求。两电子时它是
\begin{equation}
\Phi(x_1,x_2)=\frac{\phi_1(x_1)\phi_2(x_2)-\phi_2(x_1)\phi_1(x_2)}{\sqrt2}.
\end{equation}
将其代入相互作用期望值，两个模平方项各给 $J_{12}/2$，两个交叉项各给 $-K_{12}/2$，其中
\begin{equation}
K_{ij}=\iint\frac{\phi_i^*(x)\phi_j(x)\phi_j^*(x')\phi_i(x')}{|\bm r-\bm r'|}\dd x\dd x'.
\end{equation}
因此两电子相互作用为 $J_{12}-K_{12}$。一般 $N$ 电子中，对未参与算符的坐标积分时，正交性使不匹配轨道的项为零；每一对剩下直接与交换两种收缩，得到
\begin{equation}
E_{\rm HF}=\sum_i h_{ii}+\frac12\sum_{ij}(J_{ij}-K_{ij}).
\end{equation}
$i=j$ 时 $J_{ii}=K_{ii}$，故这里可把自项纳入求和而不产生自相互作用。交换不是额外的一种经典力，而是反对称态的干涉结果。若自旋轨道的自旋部分正交，交换积分的自旋积分为零，所以交换仅在同自旋通道出现。

为看其空间意义，定义一体密度矩阵 $\gamma(x,x')=\sum_i\phi_i(x)\phi_i^*(x')$，$n(x)=\gamma(x,x)$。行列式的对密度为 $n_2(x,x')=n(x)n(x')-|\gamma(x,x')|^2$。给定 $x$ 处电子后，条件密度比 $n(x')$ 少了 $|\gamma|^2/n(x)$。积分利用轨道正交性得到 $\int |\gamma(x,x')|^2\dd x'=n(x)$，因此这个交换孔的积分为 $-1$。

\section{轨道变分与 Fock 算符}
不同轨道必须相互正交，所以用矩阵乘子 $\lambda_{ij}$ 约束 $\langle\phi_i|\phi_j\rangle=\delta_{ij}$。直接项变分产生乘法算符 $\hat J_j$，交换项产生积分算符 $\hat K_j$：
\begin{align}
\hat J_j\psi(x)&=\left[\int\frac{|\phi_j(x')|^2}{|\bm r-\bm r'|}\dd x'\right]\psi(x),\\
\hat K_j\psi(x)&=\phi_j(x)\int\frac{\phi_j^*(x')\psi(x')}{|\bm r-\bm r'|}\dd x'.
\end{align}
于是 $\hat f=\hat h+\sum_j(\hat J_j-\hat K_j)$，驻值条件是 $\hat f\phi_i=\sum_j\lambda_{ji}\phi_j$。占据轨道的酉变换只使行列式改变一个整体相位，因而可在占据子空间对角化 Hermitian 乘子矩阵，得到正则轨道方程 $\hat f\phi_i=\varepsilon_i\phi_i$。交换的非局域性表现在：某点的输出依赖整条输入函数，而不只是该点的值。

乘以 $\phi_i^*$ 并积分再求和，得到 $\sum_i\varepsilon_i=\sum_i h_{ii}+\sum_{ij}(J_{ij}-K_{ij})$。这把每一对相互作用计了两次，所以
\begin{equation}
E_{\rm HF}=\sum_i\varepsilon_i-\frac12\sum_{ij}(J_{ij}-K_{ij}).
\end{equation}
对偶数电子闭壳层，每个空间轨道双占据，空间轨道形式为 $\hat f=\hat h+\sum_{j=1}^{N/2}(2\hat J_j-\hat K_j)$。因直接相互作用有两种自旋来源，而交换只有同自旋来源，系数分别是二和一。

\section{非正交基组怎样变成可解矩阵问题}
写 $\phi_i=\sum_\mu C_{\mu i}\chi_\mu$，定义重叠矩阵 $S_{\mu\nu}=\langle\chi_\mu|\chi_\nu\rangle$ 和 $F_{\mu\nu}=\langle\chi_\mu|\hat f|\chi_\nu\rangle$。投影轨道方程得到 $FC=SC\varepsilon$，而不是普通的 $FC=C\varepsilon$。

以实基函数闭壳层为例，定义 $P_{\lambda\sigma}=2\sum_{i\,\mathrm{occ}}C_{\lambda i}C_{\sigma i}$ 与
\begin{equation}
(\mu\nu|\lambda\sigma)=\iint\frac{\chi_\mu(\bm r)\chi_\nu(\bm r)\chi_\lambda(\bm r')\chi_\sigma(\bm r')}{|\bm r-\bm r'|}\dd\bm r\dd\bm r'.
\end{equation}
把轨道展开代入直接与交换项即得
\begin{equation}
F_{\mu\nu}=h_{\mu\nu}+\sum_{\lambda\sigma}P_{\lambda\sigma}[(\mu\nu|\lambda\sigma)-\tfrac12(\mu\lambda|\nu\sigma)].
\end{equation}
$F$ 依赖 $P$，而 $P$ 又依赖 $F$ 的本征向量，这就是矩阵层面的自洽循环。

若基线性无关，$S$ 正定。设 $S=UsU^\dagger$，令 $X=Us^{-1/2}U^\dagger$，则 $X^\dagger SX=I$。取 $C=XC'$，得到 Hermitian 问题 $X^\dagger FXC'=C'\varepsilon$。小重叠本征值会放大误差，应控制近线性相关，而不是直接计算一个病态的 $S^{-1}F$。每轮检查密度残差、能量变化和 $C^\dagger SC=I$，必要时混合新旧密度以减小振荡。

\section{Koopmans 结论的边界}
固定其他轨道，从占据轨道 $k$ 移去一个电子：总能量中删去 $h_{kk}$，并删去它与每个其他电子的一次相互作用。因此
\begin{equation}
E_{N-1}^{\rm frozen}-E_N=-h_{kk}-\sum_j(J_{kj}-K_{kj})=-\varepsilon_k.
\end{equation}
这给出冻结轨道 HF 电离能。真实离子会重新优化轨道，而且有超出 HF 的关联；两者都不在该推导内。对虚轨道估计电子亲和能还受束缚态是否存在影响，不能把任意轨道本征值直接称为实验能级。
